\begin{thebibliography}{99}

\bibitem[Black y Scholes(1973)]{blackscholes1973}
Black, F. y Scholes, M. (1973). The pricing of options and corporate liabilities. \textit{Journal of Political Economy}, 81(3), 637--654. https://doi.org/10.1086/260062.

\bibitem[Curran(1994)]{curran1994}
Curran, M. (1994). Valuing Asian and portfolio options by conditioning on the geometric mean price. \textit{Management Science}, 40(12), 1705--1711.

\bibitem[Guidolin y Timmermann(2003)]{guidolin2003}
Guidolin, M. y Timmermann, A. (2003). Option prices under Bayesian learning: implied volatility dynamics and predictive densities. \textit{Journal of Economic Dynamics and Control}, 27(5), 717--769. https://doi.org/10.1016/S0165-1889(01)00069-0.

\bibitem[Hastings(1970)]{hastings1970}
Hastings, W. K. (1970). Monte Carlo sampling methods using Markov chains and their applications. \textit{Biometrika}, 57(1), 97--109. https://doi.org/10.1093/biomet/57.1.97.

\bibitem[Kahalé(2020)]{kahale2020}
Kahalé, N. (2020). General multilevel Monte Carlo methods for pricing discretely monitored Asian options. \textit{European Journal of Operational Research}, 287(2), 739--748. https://doi.org/10.1016/j.ejor.2020.04.022.

\bibitem[Kemna y Vorst(1990)]{kemnavorst1990}
Kemna, A. G. Z. y Vorst, A. C. F. (1990). A pricing method for options based on average asset values. \textit{Journal of Banking \& Finance}, 14(1), 113--129. https://doi.org/10.1016/0378-4266(90)90039-5.

\bibitem[Matías González et al.(2019)]{matias2019}
Matías González, A., Martínez-Palacios, M. T. V. y Ortiz-Ramírez, A. (2019). Consumo e inversión óptimos y valuación de opciones asiáticas en un entorno estocástico con fundamentos microeconómicos y simulación Monte Carlo. \textit{Revista Mexicana de Economía y Finanzas Nueva Época}, 14(3), 397--414. https://doi.org/10.21919/remef.v14i3.408.

\bibitem[Metropolis et al.(1953)]{metropolis1953}
Metropolis, N., Rosenbluth, A. W., Rosenbluth, M. N., Teller, A. H. y Teller, E. (1953). Equation of state calculations by fast computing machines. \textit{The Journal of Chemical Physics}, 21(6), 1087--1092. https://doi.org/10.1063/1.1699114.

\bibitem[Merton(1973)]{merton1973}
Merton, R. C. (1973). Theory of rational option pricing. \textit{The Bell Journal of Economics and Management Science}, 4(1), 141--183. https://doi.org/10.2307/3003143.

\bibitem[Ortiz Ramírez y Martínez Palacios(2016)]{ortiz2016}
Ortiz Ramírez, A. y Martínez Palacios, M. T. (2016). Valuación de opciones asiáticas versus opciones europeas con tasa de interés estocástica. \textit{Contaduría y Administración}, 61(4), 629--648. https://doi.org/10.1016/j.cya.2016.06.002.

\bibitem[Ortiz Ramírez y Jiménez Preciado(2026)]{ortizjimenez2026}
Ortiz Ramírez, A. y Jiménez Preciado, A. L. (2026). Valuación de opciones asiáticas con volatilidad realizada y reversión a la media: caso CEMEXCPO. \textit{Revista Mexicana de Economía y Finanzas Nueva Época}, 21(1), e1470. https://doi.org/10.21919/remef.v21i1.1470.

\bibitem[Rombouts y Stentoft(2014)]{romboutsstentoft2014}
Rombouts, J. V. K. y Stentoft, L. (2014). Bayesian option pricing using mixed normal heteroskedasticity models. \textit{Computational Statistics \& Data Analysis}, 76, 588--605. https://doi.org/10.1016/j.csda.2013.06.023.

\end{thebibliography}
